\chapter{Practical Examples of Complex Type Questionnaires}
\label{ch:appendix-examples}

This appendix provides complete QML examples demonstrating the complex type concepts discussed in Section~\ref{sec:complex-types} and Section~\ref{sec:z3-complex-types}. These examples illustrate vector and matrix item types with their associated constraints in realistic survey scenarios.

\section{Employee Satisfaction Survey}

Consider a survey measuring employee satisfaction across multiple dimensions:

\begin{lstlisting}[language=qml, caption={Multi-dimensional satisfaction survey}, numbers=none]
qmlVersion: "1.0"
questionnaire:
  title: "Employee Satisfaction Survey"
  codeInit: |
    # Track problematic areas
    problem_areas = []

  blocks:
    - id: satisfaction_assessment
      items:
        # Item I_1: Rate 5 aspects of job satisfaction
        # Type: Vector[5], Domain: Each element in [1,10]
        - id: job_aspects
          kind: QuestionGroup
          title: "Rate your satisfaction with each aspect"
          questions:
            - "Work-life balance"
            - "Compensation and benefits"
            - "Career development"
            - "Management support"
            - "Work environment"
          input:
            control: Slider
            min: 1
            max: 10
          codeBlock: |
            # Identify problem areas (rated below 5)
            aspect_names = ["Work-life", "Compensation", "Career", "Management", "Environment"]
            problem_areas = [aspect_names[i] for i in range(5)
                           if job_aspects.outcome[i] < 5]

        # Item I_2: Detailed questions if any aspect rated below 5
        - id: followup_needed
          kind: Comment
          title: "Additional feedback needed"
          precondition:
            - predicate: any([job_aspects.outcome[j] < 5 for j in range(5)])
          codeBlock: |
            count = len(problem_areas)
            message = f"We identified {count} areas needing improvement:\n"
            for area in problem_areas:
              message += f"- {area}\n"
            print(message)

        # Item I_3: Rank problematic areas by urgency
        - id: problem_ranking
          kind: QuestionGroup
          title: "Rank these problematic areas by urgency"
          precondition:
            # Same condition as I_2
            - predicate: any([job_aspects.outcome[j] < 5 for j in range(5)])
          questions:
            - "First priority"
            - "Second priority"
            - "Third priority"
          input:
            control: Dropdown
            labels:
              1: "Work-life balance"
              2: "Compensation and benefits"
              3: "Career development"
              4: "Management support"
              5: "Work environment"
          postcondition:
            # Only allow ranking of actual problem areas
            - predicate: all([problem_ranking.outcome[i] in [j+1 for j in range(5) if job_aspects.outcome[j] < 5] for i in range(len(problem_ranking.outcome)) if problem_ranking.outcome[i] > 0])
              hint: "Please only rank areas you rated below 5"

        # Item I_4: Allocate 100 points for improvement priority
        - id: improvement_allocation
          kind: QuestionGroup
          title: "Allocate 100 points across areas for improvement"
          precondition:
            - predicate: any([job_aspects.outcome[j] < 5 for j in range(5)])
          questions:
            - "Points for Work-life balance"
            - "Points for Compensation"
            - "Points for Career development"
            - "Points for Management"
            - "Points for Work environment"
          input:
            control: Editbox
            min: 0
            max: 100
          postcondition:
            - predicate: sum(improvement_allocation.outcome) == 100
              hint: "Points must sum to exactly 100"
            - predicate: all([improvement_allocation.outcome[j] >= 0 for j in range(5)])
              hint: "All values must be non-negative"
\end{lstlisting}

This example demonstrates:
\begin{itemize}
    \item Vector items for multi-dimensional ratings using QuestionGroup
    \item Conditional logic based on vector components in preconditions
    \item Allocation constraint with sum requirement in postconditions
    \item Dynamic content generation using codeBlock
\end{itemize}

\section{Product Comparison Matrix}

A market research questionnaire comparing products:

\begin{lstlisting}[language=qml, caption={Product comparison matrix}, numbers=none]
qmlVersion: "1.0"
questionnaire:
  title: "Product Comparison Study"
  codeInit: |
    # Track best products per attribute
    best_per_attribute = []

  blocks:
    - id: comparison
      items:
        # Item I_1: Compare 4 products on 6 attributes
        # Type: Matrix[4][6], Domain: Each element in [0,100]
        - id: product_comparison
          kind: MatrixQuestion
          title: "Rate each product on these attributes (0-100)"
          rows:
            - "SmartPhone X"
            - "SmartPhone Y"
            - "SmartPhone Z"
            - "SmartPhone W"
          columns:
            - "Battery life"
            - "Screen quality"
            - "Camera"
            - "Performance"
            - "Build quality"
            - "Value for money"
          input:
            control: Editbox
            min: 0
            max: 100
          codeBlock: |
            # Calculate best product per attribute
            best_per_attribute = []
            for col in range(6):
              scores = [product_comparison.outcome[row][col] for row in range(4)]
              best_idx = scores.index(max(scores))
              best_per_attribute.append(best_idx)

        # Item I_2: Follow-up if any product scores perfect
        - id: perfect_product
          kind: Question
          title: "You rated a product perfectly! Tell us why."
          precondition:
            # Check if any product has 100 on all attributes
            - predicate: |
                any([all([product_comparison.outcome[j][k] == 100
                     for k in range(6)]) for j in range(4)])
          input:
            control: Radio
            labels:
              1: "It meets all my needs"
              2: "Best in its category"
              3: "Exceptional value"
              4: "Other reason"

        # Item I_3: Display best products per attribute
        - id: best_summary
          kind: Comment
          title: "Analysis complete"
          codeBlock: |
            # Create summary of best products
            products = ["SmartPhone X", "SmartPhone Y", "SmartPhone Z", "SmartPhone W"]
            attributes = ["Battery", "Screen", "Camera", "Performance", "Build", "Value"]
            summary = "Best products by attribute:\n"
            for i in range(6):
              best_prod = products[best_per_attribute[i]]
              score = product_comparison.outcome[best_per_attribute[i]][i]
              summary += f"- {attributes[i]}: {best_prod} (score: {score})\n"
            print(summary)
\end{lstlisting}

This demonstrates:
\begin{itemize}
    \item Matrix items for two-dimensional data using MatrixQuestion
    \item Complex preconditions involving matrix aggregations (checking for perfect scores)
    \item Derived calculations from matrix data using codeBlock
    \item Dynamic analysis and reporting of results
\end{itemize}

\clearpage
\section{Market Research Survey with Complex Dependencies}

This example demonstrates a questionnaire combining both vector and matrix types with inter-item dependencies:

\begin{lstlisting}[language=qml, caption={Questionnaire with complex type dependencies}, numbers=none]
qmlVersion: "1.0"
questionnaire:
  title: "Market Research Survey"
  codeInit: |
    # Initialize variables for tracking
    weakest_idx = -1

  blocks:
    - id: product_evaluation
      items:
        # Item I_1: Rate 4 products on 3 attributes
        # Type: Matrix[4][3] - 4 products x 3 attributes
        # Domain: Each element in [1,5]
        - id: product_ratings
          kind: MatrixQuestion
          title: "Rate each product on these attributes"
          rows:
            - "Product A"
            - "Product B"
            - "Product C"
            - "Product D"
          columns:
            - "Price"
            - "Quality"
            - "Service"
          input:
            control: Radio
            labels:
              1: "Poor"
              2: "Below Average"
              3: "Average"
              4: "Good"
              5: "Excellent"

        # Item I_2: Which product has best overall score?
        # Type: Scalar, Domain: [1,4] (product selection)
        - id: best_product
          kind: Question
          title: "Which product has the best overall score?"
          precondition:
            # At least one product must score 6+ total
            - predicate: |
                any([sum([product_ratings.outcome[i][j]
                      for j in range(3)]) >= 6
                      for i in range(4)])
          input:
            control: Radio
            labels:
              1: "Product A"
              2: "Product B"
              3: "Product C"
              4: "Product D"

        # Item I_3: Rank attributes by importance
        # Type: Vector[3] - permutation ranking
        - id: attribute_ranking
          kind: QuestionGroup
          title: "Rank these attributes by importance"
          precondition:
            # Only if chose Product A or B
            - predicate: best_product.outcome == 1 or best_product.outcome == 2
          questions:
            - "Price"
            - "Quality"
            - "Service"
          input:
            control: Dropdown
            labels:
              1: "Most important"
              2: "Second most important"
              3: "Least important"
          postcondition:
            # Ensure valid ranking (all different)
            - predicate: len(set(attribute_ranking.outcome)) == 3
              hint: "Each attribute must have a unique rank"

        # Item I_4: Rate chosen product's weakest attribute
        # Type: Scalar, Domain: [1,10]
        - id: weakest_rating
          kind: Question
          title: "Rate improvement priority for the weakest attribute"
          precondition:
            # Find and check if weakest attribute exists
            - predicate: |
                min([product_ratings.outcome[best_product.outcome-1][j]
                     for j in range(3)]) < 4
          codeBlock: |
            # Identify and display the weakest attribute
            prod_idx = best_product.outcome - 1
            scores = [product_ratings.outcome[prod_idx][j] for j in range(3)]
            weakest_idx = scores.index(min(scores))
            attr_names = ["Price", "Quality", "Service"]
            print(f"The weakest attribute is: {attr_names[weakest_idx]} (score: {scores[weakest_idx]})")
          input:
            control: Slider
            min: 1
            max: 10

        # Item I_5: Calculate weighted satisfaction scores
        # Type: Vector[3] - weighted scores per attribute
        - id: weighted_scores
          kind: Comment
          title: "Weighted satisfaction analysis completed"
          codeBlock: |
            # Calculate weighted scores based on importance ranking
            prod_idx = best_product.outcome - 1
            weighted = []
            for j in range(3):
              # Higher weight for more important attributes
              weight = 4 - attribute_ranking.outcome[j]
              score = product_ratings.outcome[prod_idx][j] * weight
              weighted.append(score)
            # Display results using built-in print function
            attr_names = ["Price", "Quality", "Service"]
            result = "Weighted scores:\n"
            for j in range(3):
              result += f"- {attr_names[j]}: {weighted[j]}\n"
            print(result)
\end{lstlisting}

This example demonstrates:
\begin{itemize}
    \item Combining matrix ratings with scalar selection and vector ranking
    \item Preconditions that reference matrix elements using dynamic indices (e.g., \texttt{product\_ratings.outcome[best\_product.outcome-1][j]})
    \item Postconditions enforcing permutation constraints on rankings
    \item Code blocks performing weighted calculations across complex types
    \item Multi-level dependencies: $I_2$ depends on $I_1$, $I_3$ depends on $I_2$, $I_4$ depends on both $I_1$ and $I_2$, and $I_5$ depends on $I_1$, $I_2$, and $I_3$
\end{itemize}
